% !TeX program = xelatex
% ============================================================
% 数学建模论文 LaTeX 模板（国赛 CUMCM 格式）
% ------------------------------------------------------------
% 版式依据：
%   1. 本地历年获奖论文版式惯例（一级标题"一、问题重述"等汉字编号、
%      二级"1.1"阿拉伯编号、关键词行"关键词: A；B；C"、
%      参考文献 [1]作者.题名[J].期刊，年(期):页码.、每问"小结"）
%   2. 参考 github.com/Mr-Potato-123/cumcm-paper-hand-skill 的
%      templates/论文模板.tex（宏包骨架、页面设置、算法/代码环境）
% 编译方式：XeLaTeX（中文支持最好）或 LuaLaTeX；编译两次处理交叉引用。
% 依赖宏包：ctex, amsmath, amssymb, booktabs, graphicx, multirow,
%           float, caption, algorithm, algorithmic, listings
% 用法：复制本文件到参赛文件夹 论文/ 目录，按注释填写内容；
%       每个子问题复制"5.1 问题一模型的建立与求解"小节的结构。
% ⚠️ 修改基线（非格式项）：
%   - 删掉了原参考模板里"以上假设的合理性将在模型检验部分进一步验证"盖章句
%     （理由不需要预告）。
% ============================================================
\documentclass[11pt]{ctexart}

% ---------- 页面设置：国赛常用 A4 上下左右 2.5cm ----------
\usepackage[a4paper, top=2.5cm, bottom=2.5cm, left=2.5cm, right=2.5cm]{geometry}

\usepackage{amsmath, amssymb}   % 数学
\usepackage{xcolor}             % 颜色（listings 高亮需要）
\usepackage{booktabs}           % 三线表
\usepackage{graphicx}           % 插图
\usepackage{multirow}           % 表格合并
\usepackage{float}              % [H] 图表位置
\usepackage{caption}            % 图表标题
\captionsetup{font=small, labelfont=bf}
\usepackage{url}

% ---------- 算法环境 ----------
\usepackage{algorithm}
\usepackage{algorithmic}
\renewcommand{\algorithmicrequire}{\textbf{输入:}}
\renewcommand{\algorithmicensure}{\textbf{输出:}}

% ---------- 代码环境（附录用） ----------
\usepackage{listings}
\lstset{basicstyle=\ttfamily\small, breaklines=true, frame=single,
        language=Python, keywordstyle=\color{blue}\bfseries,
        commentstyle=\color{gray}}

% ---------- 超链接（可选，编译用 xelatex 时不要跟 hyperref 冲突） ----------
% \usepackage{hyperref}

% ---------- 题目与作者（参赛信息：赛前填好） ----------
\title{\textbf{论文题目：研究对象与方法}}
\author{参赛队员1 \quad 参赛队员2 \quad 参赛队员3}
\date{}

\begin{document}
\maketitle

% ============================================================
% 摘要（第一页，2026 新规：电子版第一页=摘要专用页，≤1 页，最后写！）
% 每问一句：针对…问题，首先…，建立…模型，采用…算法求解，得到…结果；
%           通过…检验（直接写真实数值/百分比），表明模型…
% ⚠️ 每问叙述重心不同、句式各异，禁止三问共用同一句式只换数字（AI指纹）。
% ============================================================
\begin{abstract}
【第一段：背景 + 总体思路（本文针对…建立了…体系，得到…）】
针对问题一：【……】。
针对问题二：【……】。
针对问题三：【……】。
本文模型在【指标】上达到【数值】，可为【场景】提供参考。

\noindent\textbf{关键词:} 【词1; 词2; 词3; 词4】% 用分号分隔，3-5个
\end{abstract}

% ============================================================
% 一、问题重述（用自己的话压缩；逐问整理成明确的数学任务）
%    帮读者识别：研究对象 / 已知条件 / 核心矛盾 / 每问任务
% ============================================================
\section{问题重述}% 自动编号 "1 问题重述"（2026 红线：层次标题 1/1.1/1.1.1；正文不设目录）
\subsection{问题背景}
【用自己的话压缩题目背景，不照抄原文。】

\subsection{问题提出}
\begin{itemize}
  \item \textbf{问题一：}【整理为明确的数学任务：研究对象、已知条件、任务】。
  \item \textbf{问题二：}【同上，并说明与问题一的关系】。
  \item \textbf{问题三：}【同上】。
\end{itemize}

% ============================================================
% 二、问题分析（逐问：题目要求→关键难点→数据/变量特征→总体建模思路
%                →求解和检验方法；写"为什么该模型适合本题"，不提前展开公式）
% ============================================================
\section{问题分析}
\subsection{对问题一的分析}
【……】
\subsection{对问题二的分析}
【……注意承接问题一：是否使用问题一的模型或其输出。】
\subsection{对问题三的分析}
【……】

% ============================================================
% 三、模型假设（5-8条；只写题设之外的建模选择，不重述题设；
%              末尾不加"合理性将在检验中验证"盖章句）
% ============================================================
\section{模型假设}
\begin{enumerate}
  \item 【假设内容】【理由: 因为……所以可这样假设】;
  \item 【……】;
  \item 【……】.
\end{enumerate}

% ============================================================
% 四、符号说明（三线表；全文符号统一；放模型建立之前）
% ============================================================
\section{符号说明}
\begin{table}[H]
  \centering
  \begin{tabular}{clc}
    \toprule
    符号 & 含义 & 单位 \\
    \midrule
    $x_i$ & 第 $i$ 个…… & …… \\
    $\lambda$ & …… & …… \\
    \bottomrule
  \end{tabular}
  \caption{符号说明}
\end{table}

% ============================================================
% 五、模型的建立与求解（论文核心，占篇幅最大；每问一个闭环）
%    闭环节：问题目标→数据处理→变量定义→模型原理→模型建立→参数确认→
%            算法求解→结果呈现→结果检验→本问结论（小结）
% ============================================================
\section{模型的建立与求解}

\subsection{问题一模型的建立与求解}

\subsubsection{5.1.1 问题一的分析与数据处理}
【题目要求 $\rightarrow$ 难点 $\rightarrow$ 数据特征 $\rightarrow$ 思路。
数据预处理必须写五要素：原始数据有什么问题 / 用什么方法 / 关键参数 /
参数为什么这么选 / 处理后保留了什么特征（并展示效果）。】

\subsubsection{5.1.2 【模型名】的建立}
【先文字说明思路，再给公式，再解释每一项含义。公式前必须有引导文字
（"由……原理可得""考虑到……因素，引入……"）；公式后说明"能计算什么、下一步怎么用"。】
\begin{equation}
  f(\mathbf{x}) = \sum_{i=1}^{n} w_i x_i \label{eq:model1}
\end{equation}
其中 $w_i$ 表示【……】, $x_i$ 表示【……】。式 \eqref{eq:model1} 用于
计算【……】, 下一步将【……】。

% 优化问题给标准型：目标函数 / 决策变量 / 约束
\begin{equation}
\begin{aligned}
  \max\ / \min \quad & f(\mathbf{x}) \\
  \text{s.t.} \quad & g_j(\mathbf{x}) \le 0, \quad j = 1, \dots, m, \\
  & \mathbf{x} \in D .
\end{aligned}
\label{eq:std}
\end{equation}

\subsubsection{5.1.3 模型求解}
【算法/伪代码/流程图。流程图要体现算法内部真正步骤（初始化、迭代更新、
适应度计算、约束判断、终止条件），不可画"输入→建模→运行→输出"式无效流程图。】
\begin{algorithm}[H]
  \caption{【算法名】}
  \begin{algorithmic}[1]
    \REQUIRE 【输入: 数据、参数】
    \ENSURE 【输出: 最优解】
    \STATE 初始化种群/解集 $P_0$, 参数 $T_{\max}$;
    \FOR{$t = 1$ \textbf{to} $T_{\max}$}
      \STATE 【迭代更新规则: 如新解生成、适应度计算】;
      \STATE 【约束判断与修正】;
      \IF{满足终止条件}
        \STATE \textbf{break}
      \ENDIF
    \ENDFOR
    \RETURN 最优解
  \end{algorithmic}
\end{algorithm}

\subsubsection{5.1.4 结果与分析}
% 占位图: 换成你自己的图片; 文件不存在时自动跳过, 不影响编译
\IfFileExists{图1_xxx.png}{%
\begin{figure}[H]
  \centering
  \includegraphics[width=0.7\textwidth]{图1_xxx.png}
  \caption{图1 三种策略下的成本对比}
\end{figure}
}{}
【每张图/表必须在正文被引用并讨论（"看图说话"：图X表明……结论，……
数据如何支撑）；对比表中最优值加粗(\textbf{数值})；结果是否符合预期与假设；
检验证据：误差/残差/交叉验证/敏感性/多模型对比/与标准答案对照。】

\subsubsection{5.1.5 本问小结}
【用自己的话写本问机制性结论 + 与后问衔接；不再复读数字清单；
不以"综上"开头。如："问题一的容量约束在单向出海下退化为起飞载荷约束，
这一简化在问题二引入中途上下机后不再成立。"】

\subsection{问题二模型的建立与求解}
【同上结构。开头写承上启下过渡句："问题二在问题一的基础上进一步考虑了……,
因此我们在模型一的基础上做了如下改进……"】

\subsection{问题三模型的建立与求解}
【同上结构。】

% ============================================================
% 六、模型的评价与推广（优点 / 缺点（真实具体+改进思路）/ 改进方向 / 推广）
% ============================================================
\section{模型的评价与推广}
\subsection{模型优点}
\begin{itemize}
  \item 【……】
\end{itemize}
\subsection{模型缺点}
\begin{itemize}
  \item 【真实具体，如"计算复杂度高，数据量增大时求解时间超过……"】
\end{itemize}
\subsection{模型改进方向}
\begin{itemize}
  \item 【针对缺点给出可行改进思路】
\end{itemize}
\subsection{模型推广}
\begin{itemize}
  \item 【结合本题背景向外延伸，不空谈】
\end{itemize}

% ============================================================
% 七、参考文献（必须真实，逐一核实；8-15篇；覆盖模型/算法/领域/对比
%               四方面；只列正文引用，按引用顺序编号）
% ============================================================
\section{参考文献}
\begin{thebibliography}{9}
  \bibitem{ref1} 作者. 题名[J]. 期刊, 年, 卷(期): 页码.
  \bibitem{ref2} 作者. 题名[D]. 城市: 学校, 年.
\end{thebibliography}

% ============================================================
% 八、附录（评阅人核查真实性的地方：推导与正文一致、代码可运行、表格同源；
%          每条附录正文注明"详见附录X"，与正文一一对应，不凑数）
% ============================================================
\appendix
\section{详细推导}
【冗长推导放这里，正文注明"详细推导见附录一"。】

\section{程序代码}
\begin{lstlisting}[language=Python]
# 核心代码（删除无关调试代码，加必要注释；能原样跑出论文结果）
\end{lstlisting}

\end{document}
